\chapter{Discussion} \label{cha:discussion}

This thesis set out to investigate whether structural fault tolerance---achieved
through deliberate architectural separation and cross-modal arbitration rather
than through hardening individual models---can improve the robustness of a
heterogeneous financial AI system under adversarial signal injection. The
preceding chapter presented the empirical results of evaluating the Robust
Trinity architecture on AAPL daily equity data across clean and adversarial
conditions. This chapter interprets those results and situates them within the
broader literature. Section~\ref{sec:disc-prior-work} compares the findings
with prior work along five axes identified in the Extended Background.
Section~\ref{sec:disc-rq} provides explicit answers to the main research
question and its three sub-questions.
Section~\ref{sec:disc-limitations} discusses the limitations of the study.
Section~\ref{sec:disc-validity} addresses validity and reliability.
Section~\ref{sec:disc-ethics} considers ethical and societal implications.
Section~\ref{sec:disc-tools} describes the use of IT and AI tools throughout
the project. Section~\ref{sec:disc-future} outlines directions for future
research, and Section~\ref{sec:disc-concluding} provides a brief concluding
summary.

%% ====================================================================
\section{Comparison with Prior Work}
\label{sec:disc-prior-work}
%% ====================================================================

The Extended Background (Chapter~\ref{cha:extendBackground}) identified four
gaps in the literature: the absence of cross-modal scope in existing robustness
mechanisms, the predominant focus on training-time rather than inference-time
defences, the rarity of economically grounded evaluation, and the lack of
concrete architectural mechanisms behind governance-level safety aspirations.
This section examines how Trinity's empirical results relate to the five
principal lines of prior work surveyed there, and discusses where the present
contribution is original and where it confirms, extends, or diverges from
existing findings.

%% --------------------------------------------------------------------
\subsection{Tight-Fusion Architectures}
\label{sec:disc-tight-fusion}

The most direct point of comparison is with tight-fusion systems that embed
LLM-derived signals directly into RL trading agents.
\cite{benhenda2025finrl} demonstrate that FinRL-DeepSeek, which integrates
sentiment and risk scores from a DeepSeek LLM into a risk-sensitive RL policy,
achieves performance gains on Nasdaq-100 equities under benign conditions.
Trinity takes the opposite architectural stance: the Analyst (GPT-5) and the
Executor (PPO) share \emph{no} features, no intermediate representations, and
no gradient paths. The Analyst observes only text headlines; the Executor
observes only 14~z-normalised numeric features over a 10-day lookback window.
The C-Gate arbitrates their outputs at inference time without coupling them
during training.

The adversarial results provide direct evidence for the value of this
separation. Under analyst poisoning at all five corruption rates
(Table~\ref{tab:poison-maxdd}), the Executor-Only configuration maintains a
constant maximum drawdown of $15.28\%$, entirely unaffected by the corrupted
semantic channel. Had the Executor consumed LLM-derived features---as in a
FinRL-DeepSeek-style architecture---analyst poisoning would have propagated
into the numeric policy's action selection. The channel independence confirmed
in Chapter~\ref{cha:results} is a structural consequence of the architectural
separation, not an empirical coincidence: because the two agents share no
inputs, compromise of one channel is mathematically prevented from affecting
the other.

This does not imply that tight fusion is always inferior. Under benign
conditions, FinRL-DeepSeek's direct feature integration plausibly captures
cross-modal correlations that Trinity's separated channels cannot exploit,
which may explain why Trinity does not outperform simpler baselines on
return-based metrics. The trade-off is explicit: Trinity sacrifices the
potential performance gains of cross-modal feature sharing in exchange for a
\emph{containment guarantee} that localises the damage of adversarial
compromise to the attacked channel. Whether this trade-off is worthwhile
depends on the deployment context; in environments where adversarial risk is
material---such as production trading systems exposed to manipulated news
feeds---the containment property may outweigh the forgone return.

%% --------------------------------------------------------------------
\subsection{Trust-Weighted Aggregation}
\label{sec:disc-trust-weighted}

\cite{you_stablecoin_2026} apply trust-weighted aggregation to stablecoin
reserve management, where heterogeneous agents submit signals about reserve
adequacy and a central mechanism weights contributions by each agent's
observed behavioural reliability. Under adversarial stress testing---where a
fraction of agents submit deliberately misleading signals---the trust-weighted
scheme reduces peak peg deviation by 57\% relative to unweighted averaging.

The C-Gate can be understood as an inference-time, two-agent instantiation of
this same principle. Rather than weighting multiple homogeneous agents by
historical reliability, the C-Gate measures the real-time divergence
$\Delta_t = 1 - \pi_{\text{RL}}(d_{\text{LLM}} \mid s_t)$ between two
structurally heterogeneous agents and maps this divergence onto a three-regime
policy: full position when the agents agree ($\Delta \leq 0.6935$), half
position with a protective stop-loss during ambiguity, and a flat (cash)
position during conflict ($\Delta > 0.9877$). Both systems achieve drawdown or
deviation reduction under adversarial conditions---Trinity reduces mean maximum
drawdown by 44\% relative to Executor-Only
(Table~\ref{tab:clean-baselines}), and the stablecoin system reduces peak
deviation by 57\%. Both do so by treating inter-agent disagreement as a
first-class risk signal rather than averaging over it.

The key difference is domain and modality. The stablecoin system aggregates
agents operating within a single domain (decentralised finance reserves) over
homogeneous signal types. Trinity aggregates agents operating over
fundamentally different modalities---natural-language sentiment and numeric
price features---which means the divergence metric must bridge a semantic-numeric
gap rather than compare like with like. The C-Gate's design choice of using the
Executor's own policy probability as the bridge
($\pi_{\text{RL}}(d_{\text{LLM}} \mid s_t)$) is a direct response to this
cross-modal challenge: it avoids the need for a shared feature space by asking
how probable the Analyst's recommendation is under the Executor's learned
world model.

%% --------------------------------------------------------------------
\subsection{Byzantine Fault Tolerance and Ensemble Diversity}
\label{sec:disc-bft-ensemble}

A third line of comparison connects Trinity to work on Byzantine fault
tolerance (BFT) and adversarial robustness through ensemble diversity.
\cite{devadoss_byzantine_2025} argue that AI systems should be constructed from
multiple heterogeneous modules coordinated through consensus protocols,
borrowing the core insight of \cite{lamport1982byzantine} and
\cite{castro1999pbft}: that a system can tolerate up to $f$ faulty components
among $n$ if the remaining honest components can outvote or detect the
compromised ones. In the ensemble robustness literature,
\cite{chen_diversity_2024-1} demonstrate with the EADSR framework that
training ensemble members to produce explicitly differentiated predictions
reduces adversarial transfer, while \cite{pang_improving_2019} show that
adversarial vulnerability in ensembles scales with the degree of error
correlation across members.

Trinity can be interpreted as a minimal-case BFT mechanism operating with
$n = 2$ agents and $f = 1$ tolerated fault. The C-Gate does not implement a
majority-vote consensus---with only two agents, no majority exists when they
disagree---but it achieves an analogous function through a different mechanism:
when the two agents conflict, the system defaults to the safest possible
action (flat position) rather than trusting either agent. The empirical
results support this interpretation. Under analyst poisoning, where the
semantic channel is compromised, Trinity's maximum drawdown remains within
$7.15\%$--$9.98\%$ across all five corruption rates
(Table~\ref{tab:poison-maxdd}). Under executor perturbation, where the
numeric channel is corrupted, Trinity's maximum drawdown remains within
$6.23\%$--$9.19\%$ (Table~\ref{tab:perturb-maxdd}). In both cases, the
system contains the damage of a single-channel compromise without knowing
\emph{which} channel is compromised---the C-Gate responds to the fact of
disagreement, not to its source.

The connection to ensemble diversity is equally direct. The EADSR and related
work operate within a single modality: multiple image classifiers or multiple
time-series forecasters whose diversity is increased through training
objectives. Trinity generalises this principle from \emph{intra-modal} ensemble
diversity to \emph{inter-modal} architectural diversity. The Analyst and
Executor are diverse not because they were trained with different
objectives on the same data, but because they operate on fundamentally
different data types with fundamentally different model classes. This
structural diversity means their failure modes are inherently uncorrelated:
a perturbation crafted to fool a text-processing LLM has no mechanism to
simultaneously fool a policy network that has never seen text. The channel
independence results in Chapter~\ref{cha:results} confirm this property
empirically, providing a cross-modal analogue to the intra-modal decorrelation
that \cite{pang_improving_2019} identify as the key driver of ensemble
robustness.

%% --------------------------------------------------------------------
\subsection{Agentic Safety Frameworks}
\label{sec:disc-agentic-safety}

The agentic safety literature provides the governance context within which
Trinity's contribution should be understood. \cite{kavathekar2025tamas}
benchmark adversarial risks across multi-agent LLM configurations in the
TAMAS framework, reporting high attack success rates under impersonation and
Byzantine-style adversarial behaviour---particularly when agents can influence
one another through shared message channels. \cite{yang_finvault_2026} evaluate
execution-grounded LLM agents in FinVault, finding persistent vulnerabilities
to prompt injection and authority bypass even when guardrails are applied.
\cite{chen_standard_2025} argue that standard benchmarks underweight safety
considerations, proposing that financial agents should be evaluated by how
safely they fail rather than solely by how well they perform.
\cite{kurshan_agentic_2025} propose the Agentic Regulator, a layered
oversight architecture with monitoring at the model level, the firm level, and
the regulatory level.

Trinity maps onto this governance hierarchy in a concrete way. The C-Gate
operates at the model level: it is an inference-time mechanism that detects
cross-modal disagreement between two specific agents and translates that
disagreement into a graded response. The Guardian's rule-based risk
constraints---position sizing, stop-losses, fail-safe to cash---operate at
the firm level: they enforce risk limits that no model output can override.
Together, these two layers provide the kind of ``layered oversight'' that
\cite{kurshan_agentic_2025} advocate, but with a concrete architectural
implementation rather than a conceptual governance prescription.

The TAMAS finding that shared message channels amplify vulnerability is
particularly relevant. Trinity's channel independence is designed to
eliminate precisely this attack vector: the Analyst and Executor share no
message channel, no feature space, and no intermediate state. The C-Gate
reads both agents' outputs but does not relay one agent's output to the
other---it is a one-way arbitration layer, not a bidirectional communication
channel. This architectural choice directly addresses the structural
vulnerability that TAMAS identifies, though Trinity does so for a two-agent
financial system rather than for the general multi-agent LLM configurations
that TAMAS evaluates.

The FinVault results reinforce a complementary point. Even when
model-level guardrails are applied, \cite{yang_finvault_2026} find that
LLM agents remain vulnerable to prompt injection and authority bypass in
financial workflows. Trinity does not claim to solve this problem at the
model level---the Analyst (GPT-5) is used as a frozen, pre-trained
component without additional hardening. Instead, Trinity's defence operates
one level higher: even if the Analyst is successfully manipulated, the
C-Gate detects the resulting disagreement with the Executor and triggers a
protective response. This is a form of \emph{architectural fault
containment} rather than model-level robustness, and it complements rather
than replaces the guardrail-based defences that FinVault evaluates.

%% --------------------------------------------------------------------
\subsection{Multi-Agent Adversarial Coordination}
\label{sec:disc-multi-agent-adversarial}

Two recent studies apply multi-agent or adversarial approaches to financial
time-series problems. \cite{qiao_multi-agent_2025} propose a multi-agent
adversarial architecture for time-series forecasting (MAA-TSF), in which
heterogeneous agent populations outperform single models under volatile market
conditions. \cite{xia_bayesian_2026} introduce Bayesian robust trading with
adversarial synthetic data generators that stress-test RL trading policies
under worst-case macroeconomic regimes.

Trinity's results confirm the pattern identified by both studies: multi-agent
architectures provide measurable robustness benefits under adversarial or
volatile conditions compared to single-agent alternatives. The 44\% reduction
in maximum drawdown that Trinity achieves relative to Executor-Only under
clean conditions (Table~\ref{tab:clean-baselines}), and the tight drawdown
band maintained across all ten adversarial intensities
(Section~\ref{sec:adversarial}), are consistent with the general finding that
heterogeneous agent populations handle distributional stress better than
isolated models.

However, Trinity diverges from both studies in an important respect. MAA-TSF
and the Bayesian approach both demonstrate that their multi-agent architectures
improve \emph{return} metrics under stress. Trinity does not. None of the
active trading configurations---Trinity, Executor-Only, or Analyst-Only---achieves
a positive mean Sharpe ratio on the test period
(Table~\ref{tab:clean-baselines}), and the C-Gate ablation reveals that
Trinity-no-CGate consistently outperforms Trinity on Sharpe in every evaluation
condition (Table~\ref{tab:ablation-sharpe}). The C-Gate's value lies
exclusively in \emph{risk containment}, not in return enhancement. This
distinction is important for positioning Trinity's contribution accurately:
the architecture does not make a trading system more profitable; it makes it
more predictable under adversarial stress, at a measurable cost to return.

Whether this trade-off is attractive depends on the objective function. For
a risk-constrained institutional portfolio where drawdown limits are hard
constraints---as is common under regulatory capital requirements and internal
risk mandates---a system that maintains maximum drawdown within a 3.75~pp
band ($6.23\%$--$9.98\%$) across all adversarial conditions may be preferable
to one that occasionally achieves higher returns but with wider and less
predictable drawdown exposure.

%% --------------------------------------------------------------------
\subsection{Cross-Study Positioning}
\label{sec:disc-cross-study}

Table~\ref{tab:cross-study} summarises the adversarial robustness evaluations
reported by the studies discussed above, alongside Trinity's results, to
contextualise the present contribution within the broader landscape.

\begin{table}[h]
\centering
\caption{Qualitative comparison of adversarial robustness evaluations across
  related studies. Metrics and domains differ across studies; the table is
  intended to contextualise rather than rank.}
\label{tab:cross-study}
\small
\begin{tabular}{p{2.8cm}p{2.2cm}p{2.4cm}p{2.0cm}p{3.2cm}}
\toprule
\textsc{Study} & \textsc{Domain} & \textsc{Attack Type} &
  \textsc{Primary Metric} & \textsc{Key Finding} \\
\midrule
\cite{xia_bayesian_2026} &
  RL trading, 9~instruments &
  Adversarial macro-regime generation &
  Sharpe, profitability &
  Outperforms 9~baselines under COVID-era stress \\
\addlinespace
\cite{qiao_multi-agent_2025} &
  Time-series forecasting, 19~assets &
  Adversarial multi-agent training &
  MAE, directional accuracy &
  10--70\% MAE reduction; 5--25\% directional gains \\
\addlinespace
\cite{you_stablecoin_2026} &
  Stablecoin reserves &
  Adversarial agent injection &
  Peak peg deviation &
  57\% reduction vs.\ unweighted baseline \\
\addlinespace
\cite{goldblum_adversarial_2021} &
  HFT order book &
  Cost-constrained perturbation &
  Profitability, risk &
  Measurable degradation under budget-constrained attacks \\
\addlinespace
\cite{fursov_adversarial_2021} &
  Fraud detection &
  Synthetic transaction injection &
  False-negative rate &
  Few injected transactions fool DL models \\
\addlinespace
\cite{kavathekar2025tamas} &
  Multi-agent LLMs &
  Impersonation, Byzantine &
  Attack success rate &
  High ASR across 6~attack types and 10~LLMs \\
\addlinespace
\cite{yang_finvault_2026} &
  LLM financial agents &
  Prompt injection, authority bypass &
  Attack success rate &
  Up to 50\% ASR on SOTA models; 6.7\% on most robust \\
\addlinespace
\textbf{Trinity (this thesis)} &
  Equity trading (AAPL) &
  Signal flipping, Gaussian noise &
  MaxDD band &
  6.23--9.98\% across 10~intensities; 44\% MaxDD reduction vs.\ single-agent \\
\bottomrule
\end{tabular}
\end{table}

Several observations follow from this comparison. First, no directly
comparable state-of-the-art exists for adversarial robustness in cross-modal,
multi-agent financial trading systems. The studies in
Table~\ref{tab:cross-study} span different domains (equity trading, stablecoin
reserves, fraud detection, HFT, general forecasting), different attack models
(learned adversarial generators, random perturbations, prompt injection,
transaction injection), and different primary metrics (Sharpe ratio, MAE,
peg deviation, attack success rate, maximum drawdown). These differences
preclude direct numerical ranking and reflect the fragmented state of the
field: adversarial robustness in financial AI is evaluated through
domain-specific lenses with no shared benchmark or common evaluation protocol.

Second, the studies that are closest in spirit to Trinity---Xia et al.\ and
Qiao et al., both of which employ multi-agent or adversarial architectures for
financial decision-making---report improvements in \emph{return} and
\emph{accuracy} metrics under stress, whereas Trinity's measurable benefit is
exclusively in \emph{risk containment} (maximum drawdown stability). This
divergence underscores a fundamental design choice: Trinity prioritises
predictable worst-case behaviour over expected-case performance, a trade-off
that is explicit in the C-Gate's conflict-regime fail-safe to cash.

Third, the safety benchmarks (TAMAS, FinVault) operate at a different level
of abstraction: they measure attack success rates against model-level defences,
whereas Trinity measures the economic consequences of attacks that
\emph{succeed} at the model level but are \emph{contained} at the
architectural level. These perspectives are complementary rather than
competing---a system could incorporate both model-level guardrails (as
evaluated by FinVault) and architectural fault containment (as evaluated here)
within a layered defence.

The absence of a directly comparable baseline is itself informative. It
confirms the gap identified in the Extended Background
(Chapter~\ref{cha:extendBackground}): inference-time cross-modal arbitration
for adversarial robustness in financial trading has not been previously
evaluated, and the contribution of this thesis occupies that gap.

%% ====================================================================
\section{Answering the Research Questions}
\label{sec:disc-rq}
%% ====================================================================

This section provides explicit answers to the main research question and each
sub-research question, with references to the supporting empirical evidence.

\paragraph{SRQ1: How should semantic, numeric, and symbolic components be
decomposed and coordinated to reduce correlated failure?}

The results confirm that Trinity's channel independence---the Analyst observes
only text headlines, the Executor observes only numeric features, and neither
agent has access to the other's inputs---effectively eliminates correlated
failure across channels. Under analyst poisoning, the Executor-Only
configuration maintains constant performance metrics across all five corruption
rates ($\text{MaxDD} = 15.28\%$; Table~\ref{tab:poison-maxdd}), because the
corrupted semantic signals never enter the numeric channel. Symmetrically,
under executor perturbation, the Analyst-Only configuration maintains constant
performance ($\text{MaxDD} = 9.12\%$; Table~\ref{tab:perturb-maxdd}), because
the injected Gaussian noise in the observation vector does not propagate to
the semantic channel. This bidirectional isolation is a direct consequence of
the architectural decomposition described in Section~\ref{sec:architecture}:
by enforcing zero shared features between agents, the design ensures that
compromise of one reasoning channel cannot automatically compromise the other.

The coordination role is played by the C-Gate, which reads both agents' outputs
without coupling their inputs. The three-regime policy---full position on
agreement, reduced position on ambiguity, flat on conflict---provides a graded
response mechanism that preserves the containment benefit of channel
independence while still enabling the agents to contribute to a joint trading
decision. The answer to SRQ1 is therefore that decomposition should enforce
strict input separation across modalities, and coordination should operate
exclusively at the output level through an arbitration mechanism that treats
inter-agent disagreement as a risk signal.

\paragraph{SRQ2: How can cross-modal divergence be defined and calibrated
as a detection mechanism?}

The divergence metric $\Delta_t = 1 - \pi_{\text{RL}}(d_{\text{LLM}} \mid
s_t)$ was calibrated successfully at temperature $T = 1.0$ using raw softmax
probabilities from the Executor's policy network. The calibration procedure,
applied on the validation split (452~trading days), produced data-driven
thresholds $\tau_{\text{low}} = 0.6935$ (30th percentile) and
$\tau_{\text{high}} = 0.9877$ (80th percentile), yielding a regime
distribution of approximately 30\% agreement, 50\% ambiguity, and 20\%
conflict (Table~\ref{tab:calibration}). This distribution provides meaningful
separation: the three regimes correspond to qualitatively different levels of
inter-agent consistency, and the per-seed regime variation on the test split
(Table~\ref{tab:regime-test}) demonstrates that the metric captures genuine
differences in agent alignment rather than producing degenerate output.

The importance of temperature selection was confirmed by a negative result.
An earlier calibration at $T = 0.05$ produced degenerate thresholds
($\tau_{\text{low}} = 0.9999$, $\tau_{\text{high}} = 1.0000$) because the
sharpened softmax collapsed the policy distribution to near-one-hot vectors,
pushing $\Delta$ to a bimodal $\{0, 1\}$ distribution with no continuous
middle range. Switching to $T = 1.0$ restored a continuous $\Delta$
distribution and improved mean test Sharpe from $-0.618$ to $-0.270$
(Section~\ref{sec:calibration}). This result establishes that temperature
is a critical hyperparameter for the divergence metric's discriminative
power and must be set to preserve the full probability distribution rather
than sharpening it.

The answer to SRQ2 is that cross-modal divergence can be operationalised as
the complement of the Executor's policy probability assigned to the Analyst's
recommended action, calibrated with data-driven percentile thresholds on
validation data at a temperature that preserves the policy distribution's
continuous range. The metric successfully distinguishes varying levels of
agent disagreement and supports a graded response policy.

\paragraph{SRQ3: Does a structurally heterogeneous architecture with
cross-modal arbitration provide measurable robustness improvements compared
to single-agent and tightly fused alternatives under simulated adversarial
conditions?}

The answer is affirmative for risk containment and negative for return. Under
clean conditions, Trinity achieves the lowest mean maximum drawdown at
$8.60\%$, representing a 44\% reduction relative to the Executor-Only baseline
($15.28\%$; Table~\ref{tab:clean-baselines}). Under adversarial pressure,
Trinity's mean maximum drawdown remains within $6.23\%$--$9.98\%$ across all
ten attack intensities---a combined range of approximately 3.75~percentage
points---compared to $11.28\%$--$15.76\%$ for Executor-Only and
$4.23\%$--$12.18\%$ for Analyst-Only (Section~\ref{sec:adversarial}). The
C-Gate ablation demonstrates that the arbitration mechanism contributes a
consistent $0.45$--$0.63$~pp reduction in maximum drawdown relative to
Trinity-no-CGate across all evaluation conditions
(Table~\ref{tab:ablation}).

However, no active trading configuration achieves a positive mean Sharpe ratio
except Trinity-no-CGate ($+0.238$), and the C-Gate consistently reduces
Sharpe relative to the ablated variant in every condition
(Table~\ref{tab:ablation-sharpe}). The AAPL buy-and-hold benchmark achieves a
Sharpe of approximately $0.76$ over the same test period, outperforming all
active systems on risk-adjusted return. The robustness improvement is therefore
specific to drawdown containment: the architecture makes the system's worst-case
losses more predictable and bounded, but does not improve profitability.

\paragraph{Main RQ: How can structural fault tolerance be designed and
evaluated in heterogeneous financial AI systems to improve robustness under
adversarial signal injection?}

The results demonstrate that structural fault tolerance can be designed through
three interlocking mechanisms: (1)~channel independence, which prevents
adversarial compromise from propagating across modalities; (2)~a cross-modal
divergence metric that detects inter-agent disagreement in real time; and
(3)~a graded response policy that reduces exposure proportionally to the
detected disagreement. The evaluation framework---multi-seed assessment across
multiple adversarial attack types and intensities, measured through economic
risk metrics (maximum drawdown, Sharpe ratio, Sortino ratio)---provides a
replicable methodology for quantifying the robustness improvement.

The principal finding is that the value of structural fault tolerance lies in
risk containment rather than return enhancement. The architecture achieves
tighter and more stable drawdown bounds across adversarial conditions than
any single-agent or ablated alternative, but it does so at a cost to
return-based performance. This pattern is consistent with established
principles in fault-tolerant systems design, where redundancy and fail-safe
mechanisms improve reliability at the expense of throughput.

%% ====================================================================
\section{Limitations}
\label{sec:disc-limitations}
%% ====================================================================

The findings reported in this thesis are subject to several limitations that
constrain their generalisability, ecological validity, and the strength of
the conclusions that can be drawn.

\paragraph{Single asset and single market regime.}
All experiments were conducted on a single equity ticker (AAPL) over a
specific historical period. The training set spans approximately 756~trading
days (2021--2023), and the test set covers the second half of 2024---a period
characterised by a broadly bullish market environment for large-cap US
technology equities. The architecture's robustness properties may not
generalise to other asset classes (fixed income, commodities, foreign
exchange), other market microstructures (order-book-driven venues, over-the-counter
markets), or other market regimes (prolonged bear markets, range-bound
environments, periods of extreme volatility). The fact that none of the active
trading configurations outperforms buy-and-hold on Sharpe during the bullish
test period suggests that the trained policies may be systematically
miscalibrated to the prevailing regime, making it difficult to disentangle
architectural robustness from regime-dependent underperformance.

\paragraph{Small training dataset.}
The training set of 756~daily bars is small by the standards of deep RL, where
agents typically require millions of transitions to converge reliably. This
data limitation manifests in two ways. First, the PPO Executor exhibits
substantial cross-seed variance: individual-seed Trinity Sharpe ratios range
from $+1.240$ (seed~123) to $-2.079$ (seed~6789), a spread that indicates
high sensitivity to the stochastic elements of training (weight initialisation,
minibatch sampling). Second, the DQN executor trained during this project
similarly shows overfitting, with a validation-to-test Sharpe divergence
(mean val Sharpe $-0.115$ vs.\ mean test Sharpe $-0.746$) that further
confirms insufficient data for stable policy learning. The multi-seed
evaluation design mitigates but does not eliminate this problem: the reported
means are computed over only four seeds, which is too few for precise
confidence intervals.

\paragraph{Simulated adversarial attacks.}
The two attack types evaluated---random signal flipping and additive Gaussian
noise---are basic by the standards of adversarial machine learning. They model
an attacker who corrupts signals uniformly at random, without knowledge of the
system's architecture, thresholds, or current state. Real adversaries are
strategic: they adapt their attacks based on observed system behaviour, target
the most impactful signals, and may craft perturbations that specifically evade
detection mechanisms such as the C-Gate. The NIST taxonomy
\cite{vassilev2023nist} distinguishes between white-box, black-box, and
adaptive adversaries; the attacks evaluated here fall into the weakest
(random, untargeted) category. Trinity's robustness under these attacks is a
necessary but not sufficient condition for robustness under more sophisticated
threat models.

\paragraph{No live deployment.}
All evaluations were conducted in a simulation environment that assumes
zero-slippage execution, instantaneous order fills at the daily close price,
no transaction costs beyond those implicit in the daily return calculation,
and no market impact. In live deployment, each of these assumptions breaks
down to varying degrees depending on the traded instrument's liquidity,
the order size relative to average daily volume, and the latency of the
execution infrastructure. The gap between simulated and live performance is
well documented in quantitative finance, and the results reported here should
be understood as upper bounds on the performance achievable in practice.

\paragraph{Pre-trained frozen agents.}
Both the Analyst (GPT-5) and the Executor (PPO) are pre-trained and frozen:
they do not adapt to new data during evaluation. This design choice was
deliberate---freezing the agents ensures reproducibility and isolates the
C-Gate's contribution from confounds introduced by online learning dynamics---but
it also means the architecture cannot adapt to distributional shift in
real time. A production system operating over months or years would encounter
evolving market regimes, changing news vocabulary, and shifting cross-asset
correlations that the frozen agents cannot track. The absence of online
learning is a significant limitation for any claim about long-term deployed
robustness.

\paragraph{Analyst determinism and LLM stochasticity.}
The Analyst's signals were generated by a single GPT-5 run per headline,
with the resulting buy/hold/sell labels treated as deterministic inputs
throughout all experiments. In practice, LLM outputs are stochastic: the same
prompt can produce different completions across runs, and the degree of this
variability depends on the model's temperature setting, the decoding strategy,
and the prompt structure. By fixing a single realisation of the Analyst's
signals, the experimental design eliminates one source of variance but also
prevents any assessment of the Analyst's own reliability or the sensitivity
of the C-Gate's regime assignments to LLM output variability. Quantifying
this source of uncertainty would require multiple independent Analyst runs
per headline and propagating the resulting signal distributions through the
full evaluation pipeline.

\paragraph{Two-agent C-Gate.}
The C-Gate arbitrates between exactly two agents (Analyst and Executor). With
$n = 2$, the system can detect that disagreement exists but cannot
disambiguate which agent is compromised. In the classical BFT formulation, a
system tolerating $f$ Byzantine faults requires $n \geq 3f + 1$ participants
\cite{lamport1982byzantine}; with $n = 2$ and $f = 1$, no majority of honest
agents exists. Trinity compensates by defaulting to the safest action (flat)
during conflict, which is effective for drawdown containment but overly
conservative when the disagreement is benign. Extending the architecture to
$n > 2$ agents with a genuine consensus mechanism would address this
limitation, at the cost of increased system complexity and the need for
additional independent reasoning channels.

%% ====================================================================
\section{Validity and Reliability}
\label{sec:disc-validity}
%% ====================================================================

\paragraph{Internal validity.}
The experimental design controls for several potential confounds. The
multi-seed evaluation (four frozen seeds per configuration) mitigates the
risk that observed effects are artefacts of a single lucky or unlucky
initialisation. The ablation of the C-Gate (Trinity vs.\ Trinity-no-CGate)
isolates the arbitration mechanism's contribution by holding all other
components constant. The two adversarial attack types target different
channels independently, allowing the channel independence property to be
verified in both directions. All configurations share the same training,
validation, and test splits, the same precomputed Analyst signals, and the
same Guardian risk rules, ensuring that performance differences can be
attributed to the architectural variable under study rather than to
differences in data or auxiliary components.

One threat to internal validity is the small number of seeds. Four seeds
provide limited statistical power: the observed mean differences---such as
the $0.45$--$0.63$~pp MaxDD improvement from the C-Gate---are small relative
to the cross-seed standard deviation, and formal hypothesis tests would have
low power to reject the null. The results should therefore be interpreted as
indicative of a consistent directional effect rather than as statistically
significant in the conventional sense.

\paragraph{External validity.}
External validity is limited by the single-asset, single-period design. The
findings are specific to AAPL daily equity trading during H2~2024 and may not
transfer to other assets, frequencies, or market regimes. The simulated
adversarial attacks, while informative about the architecture's structural
properties, do not represent the full range of real-world adversarial threats.
Generalisation claims require replication across multiple assets, market
conditions, and attack models, which is beyond the scope of this thesis but
is identified as a priority for future work.

\paragraph{Reliability.}
Several design choices promote reproducibility. All random seeds are frozen
and reported. The Analyst's GPT-5 signals are precomputed and deterministic
within the experimental pipeline. Library versions (Stable-Baselines3, PyTorch,
Gymnasium) are pinned. The C-Gate calibration is fully deterministic given the
validation-set divergence distribution. Training and evaluation scripts are
version-controlled.

The principal threat to reliability is the high cross-seed variance observed
in the PPO Executor. Individual-seed Trinity Sharpe ratios range from $+1.240$
to $-2.079$---a spread of more than three Sharpe units---which means that
any single-seed result would be unreliable as a characterisation of the
architecture's performance. The multi-seed mean ($-0.270$) is more stable but
is computed over only four observations, leaving substantial uncertainty
about the true population mean. Increasing the number of seeds would improve
reliability but is constrained by the computational cost of full-pipeline
evaluation (each seed requires separate C-Gate integration and adversarial
evaluation across multiple attack types and intensities).

\paragraph{Construct validity.}
The choice of maximum drawdown as the primary robustness metric is
well-justified for the study's risk-containment framing. MaxDD directly
measures the worst peak-to-trough loss experienced by the portfolio---a
quantity that is operationally meaningful for risk management and regulatory
compliance. However, the Sharpe ratio may underweight Trinity's benefit in
contexts where drawdown constraints are binding. A system that maintains
tighter drawdown bounds but achieves lower mean return will appear inferior
on Sharpe, even if its risk-adjusted value is higher under a drawdown-penalised
objective function. The dual reporting of both MaxDD and Sharpe throughout the
results is intended to make this trade-off transparent, but readers should be
aware that the headline Sharpe comparison---in which no active system matches
buy-and-hold---may create an impression of poor performance that does not fully
reflect the architecture's risk-containment contribution.

%% ====================================================================
\section{Ethical and Societal Implications}
\label{sec:disc-ethics}
%% ====================================================================

The development and evaluation of automated trading architectures raises
several ethical and societal considerations that extend beyond the technical
results reported in this thesis.

\paragraph{Automated trading and market stability.}
Algorithmic trading systems have been implicated in episodes of rapid market
dislocation, most notably the 2010 Flash Crash, where automated selling
cascades contributed to a temporary but severe drop in US equity prices
\cite{kirilenko2017flash}. \cite{easley2011} further analyse how flow toxicity
arising from automated trading can destabilise market microstructure.
\cite{doi:10.1177/03063127211048515} examine systemic failures in algorithmic
trading more broadly, arguing that the tight coupling of automated systems
can amplify rather than contain market disruptions. While Trinity is evaluated
on daily-frequency decisions with a single ticker and minimal market impact,
the architectural principles it embodies---multi-agent coordination, automated
arbitration, fail-safe responses---would raise legitimate concerns about
systemic risk if deployed at scale across multiple assets and higher
frequencies. The fail-safe mechanism that moves Trinity to a flat position
during conflict could, if adopted by many similar systems simultaneously,
contribute to coordinated withdrawal of liquidity during periods of market
stress.

\paragraph{LLM hallucination and opacity.}
The Analyst component relies on GPT-5, a large language model whose internal
reasoning is not transparent to the system operator. LLMs are known to produce
confident but factually incorrect outputs---a phenomenon characterised as
``stochastic parrots'' by \cite{10.1145/3442188.3445922}---and the financial
domain is not exempt from this risk \cite{gehrmann2025risks}. Trinity's
architecture partially mitigates this concern by treating the Analyst's output
as one of two independent signals rather than as a trusted input, and the
C-Gate's ability to override the Analyst during conflict provides a structural
safeguard. However, the system cannot distinguish between an Analyst signal
that is wrong because of adversarial manipulation and one that is wrong because
of a hallucination. From an ethical standpoint, deploying LLM-derived trading
signals without the ability to audit the reasoning chain raises questions about
accountability when trades result in losses.

\paragraph{Risk of false confidence.}
Reporting that Trinity achieves ``robust'' performance under adversarial
conditions carries a risk of engendering false confidence in the system's
safety. The adversarial attacks evaluated are basic (random flipping, Gaussian
noise), the evaluation covers a single asset over a single period, and the
architecture has not been tested against adaptive adversaries or in live
markets. Stakeholders who interpret the robustness results without accounting
for the limitations discussed in Section~\ref{sec:disc-limitations} may
overestimate the system's readiness for production deployment. Responsible
communication of these results requires foregrounding the gap between
simulated and real-world adversarial conditions.

\paragraph{Environmental cost.}
The use of GPT-5 for generating Analyst signals, even when precomputed
offline, incurs a non-trivial computational and energy cost. The broader trend
toward LLM-integrated financial systems amplifies this concern: if every
trading decision requires an LLM inference pass, the aggregate energy
consumption and carbon footprint of the financial industry's AI infrastructure
increases correspondingly. This thesis does not quantify the environmental
cost of the Analyst's inference pipeline, but acknowledges it as a factor that
should be weighed against the marginal risk-containment benefit the Analyst
provides.

\paragraph{Responsible research practice.}
This work was conducted in accordance with the principles of good research
practice as articulated by the Swedish Research Council
\cite{vr2017goodresearchpractice}. All experimental parameters, random seeds,
and evaluation procedures are documented to enable replication. The results
are reported transparently, including negative findings (no active system
beats buy-and-hold on Sharpe) and limitations (high cross-seed variance, small
dataset, basic adversarial attacks). No proprietary or personal data were
used; the market data are publicly available historical prices and the news
headlines are drawn from publicly accessible sources.

%% ====================================================================
\section{Use of IT and AI Tools}
\label{sec:disc-tools}
%% ====================================================================

This section describes the software tools and AI systems used throughout the
research project, consistent with the requirement to disclose and discuss the
role of AI tools in the work.

The core technical infrastructure consisted of Python~3.11,
Stable-Baselines3 \cite{raffin2021stable} for PPO and DQN training,
PyTorch as the underlying deep learning framework, and Gymnasium for the
reinforcement learning environment interface. Weights~\&~Biases (WandB) was
used for experiment tracking, hyperparameter sweep management, and result
visualisation across the 50-trial DQN sweep and the multi-seed evaluation
campaigns.

The Analyst component was implemented using GPT-5 accessed through Azure
OpenAI. All Analyst signals were precomputed offline: each news headline
was processed by a single GPT-5 API call that produced a structured
buy/hold/sell label with an associated confidence score. The choice to
precompute signals ensured determinism and reproducibility across
experiments, but it also means that the Analyst's behaviour is frozen to a
single model snapshot and a single prompt template. Any changes to the
underlying model (version updates, safety filter modifications) or prompt
structure could alter the signal distribution in ways that would require
recalibration of the C-Gate thresholds.

AI-assisted coding tools were used extensively during development. GitHub
Copilot provided inline code completions during the implementation of
training scripts, evaluation pipelines, and data processing utilities.
Anthropic's Claude was used as a conversational coding assistant for
architectural design discussions, debugging, LaTeX writing, and
experiment planning. These tools accelerated development by reducing
boilerplate coding time and providing rapid feedback on implementation
choices. However, all AI-generated code was reviewed, tested, and
integrated by the author; the tools served as productivity aids rather
than autonomous contributors. The use of AI coding assistants also
introduces a reflexive consideration: a thesis investigating the
robustness of AI systems was itself partially developed with the assistance
of AI systems, which underscores both the pervasiveness of AI tools in
modern software engineering and the importance of human oversight in
ensuring the correctness and reproducibility of the resulting artefacts.

The thesis document was prepared in \LaTeX{} using Overleaf as the
collaborative editing environment. Version control was managed through Git,
with all source code, configuration files, and experiment artefacts tracked
in a single repository.

%% ====================================================================
\section{Future Research}
\label{sec:disc-future}
%% ====================================================================

The limitations identified in Section~\ref{sec:disc-limitations} suggest
several directions for future work, ranging from direct extensions of the
current architecture to more fundamental research challenges.

\paragraph{Multi-asset and multi-market generalisation.}
The most immediate priority is evaluating the Trinity architecture across
multiple assets, asset classes, and market regimes. The single-asset, single-period
design of this thesis leaves open the question of whether the robustness
properties observed for AAPL in H2~2024 are architecture-dependent or
asset-dependent. A multi-asset evaluation would also enable portfolio-level
analysis, where the interactions between per-asset C-Gate decisions and
cross-asset correlation structure could reveal emergent behaviours not
observable at the single-ticker level. Extension to different market
microstructures---order-book-driven venues, decentralised finance protocols,
commodity futures---would test whether the channel independence principle
transfers to settings with different data characteristics and latency
requirements.

\paragraph{Adaptive and strategic adversaries.}
The random attacks evaluated in this thesis represent the weakest threat
model in the adversarial machine learning taxonomy. Future work should
evaluate Trinity against adaptive adversaries that observe the system's
behaviour and optimise their perturbations to maximise damage while evading
the C-Gate's detection. This could include gradient-based attacks on the
Executor's observation space (where the attacker has white-box access to the
policy network), strategically timed headline manipulation that targets
periods of high market uncertainty (where the C-Gate's ambiguity regime is
most likely), and coordinated attacks that simultaneously perturb both
channels. The NIST taxonomy \cite{vassilev2023nist} provides a structured
framework for escalating adversarial sophistication, and the TAMAS benchmark
\cite{kavathekar2025tamas} offers concrete attack scenarios that could be
adapted to the financial trading context.

\paragraph{Online learning and continual adaptation.}
The frozen-agent design of this thesis sacrifices adaptability for
reproducibility. A natural extension is to allow the Executor to update its
policy incrementally as new market data arrive, while maintaining the
C-Gate's arbitration function. This introduces several technical challenges:
the C-Gate's calibrated thresholds may drift as the Executor's policy
evolves, requiring online recalibration; the channel independence property
must be preserved during adaptation, ruling out approaches that share
features across agents; and the risk of catastrophic forgetting in the RL
agent must be managed. Continual learning methods that maintain performance
on past regimes while adapting to new ones---such as elastic weight
consolidation or progressive neural networks---could provide a starting
point, though their application to financial RL remains largely unexplored.

\paragraph{DQN executor as an alternative to PPO.}
The DQN infrastructure developed during this project provides a second
executor implementation with a discrete action space and value-based
learning. Preliminary multi-seed results show a tighter Q-value spread
across seeds ($1.167 \pm 0.022$) but similar overfitting patterns to PPO
(validation Sharpe $-0.115 \pm 1.272$ vs.\ test Sharpe $-0.746 \pm 1.656$).
A head-to-head comparison of DQN and PPO executors within the full Trinity
pipeline---including C-Gate integration and adversarial evaluation---would
determine whether the choice of RL algorithm materially affects the
architecture's robustness properties, or whether the robustness benefits are
primarily a function of the architectural design (channel independence, C-Gate
arbitration) rather than the specific executor implementation.

\paragraph{Full BFT extension with $n > 2$ agents.}
The two-agent C-Gate cannot disambiguate which agent is compromised during
conflict. Extending the architecture to three or more independently reasoning
agents---each consuming a different modality or data source---would enable
majority-vote consensus and allow the system to identify and isolate the
compromised component. For example, a third agent consuming alternative data
(satellite imagery, supply chain signals, options flow) could serve as an
independent tiebreaker. The classical BFT threshold of $n \geq 3f + 1$
\cite{castro1999pbft} suggests that tolerating one Byzantine fault requires
at least four agents, which is more demanding than the current two-agent
design but potentially achievable with the diversity of data sources
available in modern financial markets.

\paragraph{Learned C-Gate thresholds.}
The current thresholds ($\tau_{\text{low}}$, $\tau_{\text{high}}$) are set
as fixed percentiles of the validation-set $\Delta$ distribution. An
alternative approach would learn these thresholds jointly with a downstream
objective---for example, minimising maximum drawdown subject to a return
constraint---using the validation set as a tuning set. This could be
implemented as a simple grid search, a Bayesian optimisation procedure, or
a differentiable threshold layer trained end-to-end with the Guardian's
risk rules. The risk of overfitting the thresholds to the validation period
would need to be managed through appropriate regularisation or
cross-validation.

\paragraph{Integration with regulatory compliance frameworks.}
The layered oversight structure of the Agentic Regulator
\cite{kurshan_agentic_2025} suggests a path toward integrating Trinity's
C-Gate with formal regulatory compliance mechanisms. The C-Gate's regime
assignments (agreement, ambiguity, conflict) and the associated trading
actions could be logged and audited as part of a model risk management
framework, providing regulators with a transparent record of when the system
detected inter-agent disagreement and how it responded. Mapping the C-Gate's
behaviour onto regulatory requirements for model validation, stress testing,
and explainability would require collaboration with compliance practitioners
and is beyond the scope of this technical thesis, but it represents a
practical direction for bridging the gap between architectural robustness
research and production deployment.

%% ====================================================================
\section{Concluding Summary}
\label{sec:disc-concluding}
%% ====================================================================

The Robust Trinity architecture demonstrates that structural fault
tolerance---achieved through channel independence, cross-modal divergence
detection, and a graded response policy---provides measurable robustness
benefits under simulated adversarial conditions, specifically in the form of
tighter and more stable maximum drawdown bounds. These benefits come at a
cost to return-based performance: the C-Gate's conservative fail-safe policy
reduces Sharpe relative to configurations that bypass the arbitration
mechanism. The architecture's value proposition is therefore specific to
risk-constrained settings where predictable worst-case behaviour is more
important than maximising expected return. The limitations of the study---single
asset, small dataset, basic adversarial attacks, frozen agents---constrain
the strength of the conclusions but do not undermine the central finding that
architectural separation and inference-time arbitration represent a viable and
previously unexplored approach to robustness in heterogeneous financial AI.
